\section{Conclusion}
\label{sec: conclusion}

In this paper, we focus on the estimation of a hierarchical posterior with Monte Carlo integrals. We reconcile two apparently competing ideas for tracking the accuracy of the posterior estimator with a unified statistic $\hat{E}[\hat{p}]$ for calculating the average information loss due to the approximation. We show that the likelihood estimator suffers from a bias which we argue can be easily corrected, but that the posterior estimator is also biased and, though it may be corrected in principle (examples in Appendix~\ref{app: examples of posterior correction}), it is expensive. In practice, the posterior bias is smaller than the uncertainty in the posterior estimator, and to this end, we derive a precision statistic to quantify this uncertainty (similar to metrics introduced in \citet{Essick:2022ojx}; \citet{Talbot:2023pex}). We suggest that analysts should ensure their posterior estimator does not suffer from significant Monte Carlo error by calculating the error statistic in post-processing and verifying $\hat{E}[\hat{p}] \lesssim 0.2$ bits. 

To enable analysts to compute the error statistics, we also publish a \texttt{pip}-installable python package called \texttt{population-error} \citep{population_error2025github} powered by \texttt{jax} \citep{jax2018github}.  \texttt{population-error} takes as input the set of $\Nsamp$ samples from the posterior estimator $\hat{p}(\Lambda | \{d_i\})$, the population model $p(\theta|\Lambda)$ and the $\Nobs\times\NPE$ posterior samples and $\Nfound$ found injections for computing the accuracy, precision and error statistics. We discuss the calculation of these statistics in greater detail in Appendix~\ref{app: computing the statistics}.


% An additional technique for estimating the systematic uncertainty from Monte Carlo approximations was discussed in \citet{Talbot:2023pex}. The idea is to construct the estimate of the covariance matrix $\hat{C}_{\ln\mathcal{L}}(\Lambda_i, \Lambda_j)$ in the log-likelihood for the set of samples from the posterior estimator $\Lambda_i\sim\hat{p}(\Lambda)$. Then, draw \textit{log-weights} from the Gaussian process, and reweight the posterior according to the normalized weights. Repeating this many times obtains an estimate of the systematic uncertainty due to the Monte Carlo approximation.

% Future work could extend these error statistics to additionally compute the expected information loss in marginal posteriors, e.g. the information lost in the marginal posterior on the slope of the mass distribution. 
Even with current catalog sizes of $\mathcal{O}(100)$ events, posterior uncertainty is a major systematic source of error. We show in Sec.~\ref{sec: gw example} that the standard Monte Carlo approach may remain reliable in certain simple models as the number of observations increases, but only under somewhat idealized assumptions (e.g. that simple models will remain appropriate for informative data). It quickly becomes computationally intractable for weak models and large catalogs. Ultimately, a new approach to avoid Monte Carlo uncertainties entirely may become necessary. One idea jointly samples the hierarchical posterior by fitting density estimates to each single event, promising to avoid the single-event Monte Carlo integrals \citep{Rinaldi:2021bhm,Callister:2022qwb,Hussain:2024qzl,Mancarella:2025uat}. However, it will also be important to remove the Monte Carlo approximation to the selection efficiency, which may be tackled by analytic or machine learning approaches \citep{Lorenzo-Medina:2024opt, Talbot:2020oeu, Gerosa:2020pgy, Callister:2024qyq}. Another promising technique proposes avoiding approximate likelihoods entirely via likelihood-free, simulation-based inference \citep{Leyde:2023iof}.
Though the impact of uncertainty due to the approximations in these proposals are typically neglected, we argue that our approach (though based on the common Monte Carlo method for likelihood estimation) can be extended to these in the future as long as one can construct an estimate for the covariance in the log-likelihood surface. 

\section{Acknowledgments}

We gratefully acknowledge helpful conversations with Reed Essick, Will Farr, Jacob Golomb, Matthew Mould and Colm Talbot. 
%
Additionally, we thank Reed Essick for agreeing to perform the LIGO internal document review.
%
J.H. is supported by the NSF Graduate Research Fellowship under Grant No. DGE1122374.
%
J.H. and S.V. are partially supported by the NSF grant PHY-2045740.
%
This material is based upon work supported by NSF's LIGO Laboratory which is a major facility fully funded by the National Science Foundation.
%
The authors are grateful for computational resources provided by the LIGO Laboratory and supported by National Science Foundation Grants PHY-0757058 and PHY-0823459.